\section{Система профилирования}

\subsection{Архитектура профилировщика}

\subsubsection{Кодогенерируемые профилировочные хуки}

В основе системы профилирования лежит принцип разделения: метаданные метрик хранятся в JSON-манифесте
\texttt{profiling\_metrics.json}, а код реализации генерируется автоматически скриптом \texttt{autogen.py}.

Подход основан на том, что набор метрик~--- изменяемая
часть системы, тогда как механизм сбора и передачи данных стабилен.
Перенос метаданных в JSON-манифест позволяет добавлять новые метрики
без риска внести ошибки в горячий путь исполнения инструкций:

\begin{lstlisting}[ caption={Фрагмент profiling\_metrics.json}]
{
  "instructions": {
    "*":  { "metrics": ["branch_efficiency"] },
    "ld": { "metrics": ["branch_efficiency",
                        "bank_conflicts",
                        "global_mem_transactions",
                        "global_coalescing"] },
    "st": { "metrics": ["branch_efficiency",
                        "global_mem_transactions",
                        "global_coalescing"] }
  }
}
\end{lstlisting}

Ключ \texttt{"*"} задаёт метрики для всех инструкций;
специализированные ключи добавляют метрики для конкретных мнемоник.
Добавление новой метрики~--- изменение JSON и перезапуск
\texttt{autogen.py}~--- без правки C++-кода.

Специализация метрик по мнемонике (например, \texttt{bank\_conflicts}
только для \texttt{ld}) позволяет избежать бессмысленных вычислений:
применение метрики конфликтов банков к арифметической инструкции
не имеет семантического смысла и лишь увеличивало бы объём собираемых данных.

Файл \texttt{profilers.tpp} содержит \texttt{inline}-функции
вычисления метрик и подключается один раз в начале
\texttt{instructions.cpp}. Для каждой инструкции \texttt{autogen.py}
генерирует два метода: \texttt{Execute()} и \texttt{Profile()}.
В конце \texttt{Execute()} под атрибутом \texttt{[[unlikely]]}
добавляется вызов \texttt{Profile()}:

\begin{lstlisting}[language=C++, caption={Автогенерируемые Execute() и Profile() для ld}]
void ldInstruction::Execute(std::shared_ptr<WarpContext>& wc) {
    // dispatch ExecuteThread<Data> for each active lane
    if (wc->profiling_buf) [[unlikely]]
        Profile(wc, wc->pc, wc->execution_mask, *wc->profiling_buf);
    wc->pc++;
}

void ldInstruction::Profile(std::shared_ptr<WarpContext>& wc,
        uint64_t pc, uint32_t orig_mask, WarpProfilingBuffer& buf) {
    ProfilingRecord rec;
    rec.pc = pc;  rec.instr_name = "ld";
    rec.metrics.emplace_back("branch_efficiency",
        FormatFloat(ComputeBranchEfficiency(*this, orig_mask, wc)));
    rec.metrics.emplace_back("bank_conflicts",
        std::to_string(ComputeBankConflicts(*this, orig_mask, wc)));
    rec.metrics.emplace_back("global_coalescing",
        FormatFloat(ComputeGlobalCoalescing(*this, orig_mask, wc)));
    buf.push_back(std::move(rec));
}
\end{lstlisting}

Атрибут \texttt{[[unlikely]]} обеспечивает минимальные накладные расходы
на горячем пути: при выключенном профилировании предсказатель ветвлений
не входит в \texttt{Profile()}.

Подобная организация кода позволяет использовать эмулятор в двух режимах
без каких-либо изменений бинарного файла: в режиме без профилирования
стоимость проверки флага сводится к предсказанной ветви, что практически
не отличается от режима без проверки~--- разница в производительности
определяется исключительно объёмом собираемых данных.

\subsubsection{Структура \texttt{ProfilingRecord} и буферизация}

Каждый момент исполнения инструкции порождает запись:

\begin{lstlisting}[language=C++, caption={Структура ProfilingRecord}]
struct ProfilingRecord {
    uint64_t    timestamp_ns;
    uint64_t    pc;
    uint32_t    block_id;
    uint32_t    warp_id;
    uint32_t    execution_mask;
    std::string function_name;
    std::string basic_block;
    std::string instr_name;
    std::vector<std::pair<std::string, std::string>> metrics;
};
\end{lstlisting}

Стратегия буферизации: каждый варп накапливает записи локально в \\
\texttt{WarpProfilingBuffer} без синхронизации. Сброс в общий
\texttt{Profiler} под мьютексом выполняется \textit{один раз} по
завершении блока. Число синхронизаций:
$N_{\mathrm{sync}} = N_{\mathrm{blocks}} \ll N_{\mathrm{records}}$,
что минимизирует накладные расходы.

Выбор момента сброса буфера~--- по завершении блока, а не по каждой записи~---
является ключевым для производительности: при типичном ядре один блок
генерирует тысячи записей, и синхронизация на каждую запись многократно увеличивала бы константный множитель линейных накладных расходов, что неприемлемо при тысячах записей на блок.

\subsubsection{Метка времени: \texttt{steady\_clock}}

\begin{lstlisting}[language=C++, caption={Получение временной метки}]
auto ts = std::chrono::steady_clock::now();
uint64_t timestamp_ns = std::chrono::duration_cast<
    std::chrono::nanoseconds>(ts - epoch_).count();
\end{lstlisting}

Использование \texttt{steady\_clock} (монотонные часы) вместо
\texttt{system\_clock} гарантирует монотонное возрастание меток времени
независимо от корректировок NTP, что обеспечивает корректную сортировку
записей при многоблочном исполнении.

\subsection{Формальные определения метрик}

\subsubsection{\texttt{branch\_efficiency}: эффективность дивергенции ветвлений}

Пусть $M_k$~--- маска исполнения варпа на $k$-м исполнении инструкции,
$W$~--- размер варпа (32). Для инструкций без явного предикатного
операнда эффективность исполнения:
\begin{equation}
  \label{eq:branch_eff_ch7}
  \eta_k = \frac{\popcount(M_k)}{W}.
\end{equation}

Для инструкций с явным предикатом \texttt{@\%p} дополнительно
учитывается расхождение путей «исполняется» / «пропускается»
внутри активной маски:
\begin{equation}
  \eta_k = \frac{\max\bigl(\popcount(M_k \wedge P),\;
                           \popcount(M_k \wedge \neg P)\bigr)}
                {\popcount(M_k)}.
\end{equation}

Среднее по всем запускам ядра:
\begin{equation}
  \overline{\eta} = \frac{1}{K}\sum_{k=1}^{K} \eta_k,
  \quad K = \mathit{issue\_count}.
\end{equation}

\subsubsection{\texttt{bank\_conflicts}: конфликты банков разделяемой памяти}

Разделяемая память организована в $B = 32$ банка шириной
$w = 4$~байта~\cite{CUDAProgrammingGuide, AmpereWhitepaper}. Номер банка для адреса $a$:
\begin{equation}
  \mathrm{bank}(a) = \left\lfloor \frac{a}{w} \right\rfloor \bmod B.
\end{equation}

Число конфликтов при одном исполнении \texttt{ld.shared}/\texttt{st.shared}:
\begin{equation}
  \label{eq:bank_conflicts_ch7}
  C = \max_{b \in [0,B)}\bigl|\{t \in \activeset(M) \mid \mathrm{bank}(ea(t)) = b\}\bigr| - 1.
\end{equation}

$C = 0$ означает отсутствие конфликтов (все потоки обращаются в разные
банки); $C = k$ означает $k$-кратный конфликт. Вычисление имеет
сложность $O(|\activeset(M)| + B) = O(1)$ при фиксированных $W{=}32$, $B{=}32$ на исполнение.

\subsubsection{\texttt{global\_coalescing}: коэффициент слияния транзакций}

Кэш-строка глобальной памяти занимает 128~байт. Число фактических
транзакций при одном исполнении:
\begin{equation}
  \label{eq:t_actual_ch7}
  T_{\mathrm{actual}} =
    \left|\left\{\left\lfloor\frac{ea(t)}{128}\right\rfloor
    \;\middle|\; t \in \activeset(M)\right\}\right|.
\end{equation}

Теоретический минимум транзакций при последовательном расположении
$|\activeset(M)|$ элементов размером $s$~байт:
\begin{equation}
  T_{\mathrm{ideal}} =
    \left\lceil \frac{|\activeset(M)| \cdot s}{128} \right\rceil.
\end{equation}

Коэффициент слияния:
\begin{equation}
  \label{eq:coalescing_ch7}
  \kappa = \min\!\left(1,\;\frac{T_{\mathrm{ideal}}}{T_{\mathrm{actual}}}\right)
  \in (0, 1].
\end{equation}

$\kappa = 1$ соответствует идеально слитному доступу; $\kappa \ll 1$~---
сильно разрозненному (stride-доступ, scatter/gather).

\subsubsection{\texttt{write\_ops}: давление на регистровый файл}

Для каждого типа регистров (префикс \texttt{\%r}, \texttt{\%f},
\texttt{\%rd}, \ldots) накапливается суммарное число записей:
\begin{equation}
  W[\tau] = \sum_{\substack{k:\;\text{инструкция }k\\\text{пишет в регистр типа }\tau}}
  \popcount(M_k).
\end{equation}

Инструкции без регистра-назначения (\texttt{st}, \texttt{bar},
\texttt{bra}, \texttt{ret}) в сумму не входят.

\subsection{Инструмент генерации отчётов}

\subsubsection{Анализатор текстового вывода (\texttt{parser.py})}

Файл \texttt{profiling.txt} имеет иерархическую структуру:
«запуск ядра → блок → варп → инструкция». Анализатор реализует конечный
автомат:

\begin{enumerate}
  \item \texttt{await\_launch}: ожидание строки
        \texttt{"[LAUNCH] kernel\_name grid(…) block(…)"}.
  \item \texttt{await\_block}: ожидание строки
        \texttt{"[BLOCK] block\_id=…"}.
  \item \texttt{in\_block}: разбор строк записей регулярным
        выражением \texttt{\_RECORD\_RE} за O(1) на строку.
\end{enumerate}

Результат: \texttt{list[ProfilingRecord]} для передачи в агрегатор.

\subsubsection{PTX-анализатор для отчёта (\texttt{ptx\_parser.py})}

Python-реализация анализатора решает задачу построения индекса:
\begin{equation}
  \mathit{by\_func\_pc} : (\mathit{fn},\, \mathit{pc})
  \to \mathit{PtxInstr},
\end{equation}
обеспечивающего O(1)-соответствие между профилировочной записью
и PTX-строкой для аннотации исходного кода.

\subsubsection{Агрегатор (\texttt{aggregator.py}): формулы накопления}

Для каждого уникального PC:
\begin{align}
  \mathit{exec\_count} &\mathrel{+}= \popcount(M_k), \\
  \mathit{issue\_count} &\mathrel{+}= 1, \\
  \mathit{be\_sum} &\mathrel{+}= \eta_k, \\
  \mathit{bc\_total} &\mathrel{+}= C_{\mathrm{bank},k}, \\
  \mathit{kappa\_sum} &\mathrel{+}= \kappa_k.
\end{align}

Средние значения:
\begin{equation}
  \overline{\eta}_{\mathrm{branch}} = \frac{\mathit{be\_sum}}{\mathit{issue\_count}},
  \qquad
  \bar{\kappa} = \frac{\mathit{kappa\_sum}}{\mathit{issue\_count}}.
\end{equation}

\subsubsection{Рендерер HTML-отчёта (\texttt{renderer.py})}

Самодостаточный HTML-файл содержит пять секций:

\begin{enumerate}
  \item \textbf{Панель «пилюль»}: сводные метрики ядра~---
        число инструкций, потоков, среднее $\overline{\eta}_{\mathrm{branch}}$,
        общие банковые конфликты, средний $\bar{\kappa}$.
  \item \textbf{Таблица инструкций}: все уникальные PC с тепловой
        подсветкой по \texttt{exec\_count}.
  \item \textbf{Горячие точки}: топ-10 инструкций по дивергенции,
        конфликтам банков, коэффициенту слияния.
  \item \textbf{Таблица регистров}: для каждого префикса~---
        тип, объявлено слотов, накопленный \texttt{write\_ops}.
  \item \textbf{Аннотированный исходный код}: строки исходного
        \texttt{.cu}-файла с метриками инструкций, привязанных
        через \texttt{.loc}.
\end{enumerate}

Все CSS и JS встроены в HTML~--- один файл без внешних зависимостей.

Структура из пяти секций отражает уровни детализации, необходимые
разработчику на разных этапах анализа: сводные метрики позволяют
быстро оценить наличие проблем, горячие точки сужают поиск,
аннотация исходного кода указывает на конкретную строку для доработки.
Такая иерархия соответствует типичному рабочему процессу оптимизации
CUDA-ядер и позволяет перейти от «что плохо» к «где исправлять»
в рамках одного отчётного документа.

\begin{figure}[h]
  \centering
  \includegraphics[width=16cm]{profiler_html.png}
  \caption{HTML-отчёт профилировщика: панель сводных метрик для ядра \texttt{fft\_custom}.}
  \label{fig:profiler_html}
\end{figure}

\newpage
